% GENERATED FILE. Edit canonical lessons, then run npm run generate:books.
% canonical-source-hash: fnv1a64:8d004f044e591a51
% canonical-lessons: JA-C522-nanika, JA-C522-dareka, JA-C522-dokoka, JA-C522-itsudemo, JA-C522-nandemo

\chapter{Something, Someone, Somewhere}
\label{ch:ja-something-someone-somewhere}
\hlchaptermodality{522}

\begin{chapteropening}
\textbf{By the end of this chapter:} \emph{I can ask for help if anything happens.}

Complete the last lesson of chapter 522: I can ask for help if anything happens.
\end{chapteropening}

\section[\texorpdfstring{\ja{なにか} (nanika)}{なにか (nanika)}]{\ja{なにか} (nanika) — something; anything}

\label{lesson:JA-C522-nanika}



\begin{quote}
Before the new one: say the Japanese for finally (after a long wait), then the Japanese for a good deal.
\end{quote}

\subsection*{You'll want to know: \ja{なにか}}
\textbf{\ja{なにか}} — \emph{nanika} — \textquotedblleft{}something; anything\textquotedblright{}.

\begin{quote}
\textbf{\ja{もし} \ja{なにか} \ja{あったら、でんわして} \ja{ください。}} — \emph{moshi nanika attara, denwa shite kudasai} — \textquotedblleft{}If anything happens, please call me.\textquotedblright{}
\end{quote}

\begin{cousinweb}
\textbf{\ja{なにか}} (\emph{nanika}) is a native Japanese word (wago), inherited rather than borrowed from Chinese: built from \emph{nani}, \textquotedblleft{}what\textquotedblright{}, and the particle \emph{ka} (JMdict).
\end{cousinweb}

\begin{grammarlens}[title={What you've built}]
One more word for saying what happens if.
\end{grammarlens}

\subsection*{Your turn}
\begin{itemize}
\raggedright
  \item \emph{Say it:} \emph{nanika}
  \item \emph{Say it:} \emph{nanika}, once more
  \item \emph{Say it:} \emph{moshi nanika attara, denwa shite kudasai} — If anything happens, please call me.
  \item \emph{Recall:} say the Japanese for finally (after a long wait), then the Japanese for a good deal, then say \emph{nanika} again
\end{itemize}

\begin{tcolorbox}[breakable,colback=teal!4,colframe=teal!35!black,title={Before you move on}]
What does \ja{なにか} mean? (\textquotedblleft{}Something; anything\textquotedblright{}.) Say it once more.
\end{tcolorbox}
\section[\texorpdfstring{\ja{だれか} (dareka)}{だれか (dareka)}]{\ja{だれか} (dareka) — someone; anyone}

\label{lesson:JA-C522-dareka}



\begin{quote}
Before the new one: say the Japanese for a good deal, then the Japanese for something.
\end{quote}

\subsection*{You'll want to know: \ja{だれか}}
\textbf{\ja{だれか}} — \emph{dareka} — \textquotedblleft{}someone; anyone\textquotedblright{}.

\begin{quote}
\textbf{\ja{もし} \ja{だれか} \ja{きたら、おしえて} \ja{ください。}} — \emph{moshi dareka kitara, oshiete kudasai} — \textquotedblleft{}If someone comes, please tell me.\textquotedblright{}
\end{quote}

\begin{cousinweb}
\textbf{\ja{だれか}} (\emph{dareka}) is a native Japanese word (wago), inherited rather than borrowed from Chinese: built from \emph{dare}, \textquotedblleft{}who\textquotedblright{}, and the particle \emph{ka} (JMdict).
\end{cousinweb}

\begin{grammarlens}[title={What you've built}]
One more word for saying what happens if.
\end{grammarlens}

\subsection*{Your turn}
\begin{itemize}
\raggedright
  \item \emph{Say it:} \emph{dareka}
  \item \emph{Say it:} \emph{dareka}, once more
  \item \emph{Say it:} \emph{moshi dareka kitara, oshiete kudasai} — If someone comes, please tell me.
  \item \emph{Recall:} say the Japanese for a good deal, then the Japanese for something, then say \emph{dareka} again
\end{itemize}

\begin{tcolorbox}[breakable,colback=teal!4,colframe=teal!35!black,title={Before you move on}]
What does \ja{だれか} mean? (\textquotedblleft{}Someone; anyone\textquotedblright{}.) Say it once more.
\end{tcolorbox}
\section[\texorpdfstring{\ja{どこか} (dokoka)}{どこか (dokoka)}]{\ja{どこか} (dokoka) — somewhere}

\label{lesson:JA-C522-dokoka}



\begin{quote}
Before the new one: say the Japanese for something, then the Japanese for someone.
\end{quote}

\subsection*{You'll want to know: \ja{どこか}}
\textbf{\ja{どこか}} — \emph{dokoka} — \textquotedblleft{}somewhere\textquotedblright{}.

\begin{quote}
\textbf{\ja{もし} \ja{どこか} \ja{いい} \ja{みせが} \ja{あったら、おしえて} \ja{ください。}} — \emph{moshi dokoka ii mise ga attara, oshiete kudasai} — \textquotedblleft{}If there is a good shop somewhere, please tell me.\textquotedblright{}
\end{quote}

\begin{cousinweb}
\textbf{\ja{どこか}} (\emph{dokoka}) is a native Japanese word (wago), inherited rather than borrowed from Chinese, and usually written in kana: built from \emph{doko}, \textquotedblleft{}where\textquotedblright{}, and the particle \emph{ka} (JMdict).
\end{cousinweb}

\begin{grammarlens}[title={What you've built}]
One more word for saying what happens if.
\end{grammarlens}

\subsection*{Your turn}
\begin{itemize}
\raggedright
  \item \emph{Say it:} \emph{dokoka}
  \item \emph{Say it:} \emph{dokoka}, once more
  \item \emph{Say it:} \emph{moshi dokoka ii mise ga attara, oshiete kudasai} — If there is a good shop somewhere, please tell me.
  \item \emph{Recall:} say the Japanese for something, then the Japanese for someone, then say \emph{dokoka} again
\end{itemize}

\begin{tcolorbox}[breakable,colback=teal!4,colframe=teal!35!black,title={Before you move on}]
What does \ja{どこか} mean? (\textquotedblleft{}Somewhere\textquotedblright{}.) Say it once more.
\end{tcolorbox}
\section[\texorpdfstring{\ja{いつでも} (itsudemo)}{いつでも (itsudemo)}]{\ja{いつでも} (itsudemo) — at any time}

\label{lesson:JA-C522-itsudemo}



\begin{quote}
Before the new one: say the Japanese for someone, then the Japanese for somewhere.
\end{quote}

\subsection*{You'll want to know: \ja{いつでも}}
\textbf{\ja{いつでも}} — \emph{itsudemo} — \textquotedblleft{}at any time\textquotedblright{}.

\begin{quote}
\textbf{\ja{もし} \ja{こまったら、いつでも} \ja{でんわして} \ja{ください。}} — \emph{moshi komattara, itsudemo denwa shite kudasai} — \textquotedblleft{}If you are in trouble, call me any time.\textquotedblright{}
\end{quote}

\begin{cousinweb}
\textbf{\ja{いつでも}} (\emph{itsudemo}) is a native Japanese word (wago), inherited rather than borrowed from Chinese, and usually written in kana: built from \emph{itsu}, \textquotedblleft{}when\textquotedblright{}, and \emph{demo}, \textquotedblleft{}even\textquotedblright{} (JMdict).
\end{cousinweb}

\begin{grammarlens}[title={What you've built}]
One more word for saying what happens if.
\end{grammarlens}

\subsection*{Your turn}
\begin{itemize}
\raggedright
  \item \emph{Say it:} \emph{itsudemo}
  \item \emph{Say it:} \emph{itsudemo}, once more
  \item \emph{Say it:} \emph{moshi komattara, itsudemo denwa shite kudasai} — If you are in trouble, call me any time.
  \item \emph{Recall:} say the Japanese for someone, then the Japanese for somewhere, then say \emph{itsudemo} again
\end{itemize}

\begin{tcolorbox}[breakable,colback=teal!4,colframe=teal!35!black,title={Before you move on}]
What does \ja{いつでも} mean? (\textquotedblleft{}At any time\textquotedblright{}.) Say it once more.
\end{tcolorbox}
\section[\texorpdfstring{\ja{なんでも} (nandemo)}{なんでも (nandemo)}]{\ja{なんでも} (nandemo) — anything; whatever}

\label{lesson:JA-C522-nandemo}



\begin{quote}
Before the new one: say the Japanese for somewhere, then the Japanese for at any time.
\end{quote}

\subsection*{You'll want to know: \ja{なんでも}}
\textbf{\ja{なんでも}} — \emph{nandemo} — \textquotedblleft{}anything; whatever\textquotedblright{}.

\begin{quote}
\textbf{\ja{もし} \ja{しつもんが} \ja{あったら、なんでも} \ja{きいて} \ja{ください。}} — \emph{moshi shitsumon ga attara, nandemo kiite kudasai} — \textquotedblleft{}If you have questions, ask me anything.\textquotedblright{}
\end{quote}

\begin{cousinweb}
\textbf{\ja{なんでも}} (\emph{nandemo}) is a native Japanese word (wago), inherited rather than borrowed from Chinese, and usually written in kana: built from \emph{nan}, \textquotedblleft{}what\textquotedblright{}, and \emph{demo}, \textquotedblleft{}even\textquotedblright{} (JMdict).
\end{cousinweb}

\begin{grammarlens}[title={What you've built}]
One more word for saying what happens if.
\end{grammarlens}

\subsection*{Your turn}
\begin{itemize}
\raggedright
  \item \emph{Say it:} \emph{nandemo}
  \item \emph{Say it:} \emph{nandemo}, once more
  \item \emph{Say it:} \emph{moshi shitsumon ga attara, nandemo kiite kudasai} — If you have questions, ask me anything.
  \item \emph{Recall:} say the Japanese for somewhere, then the Japanese for at any time, then say \emph{nandemo} again
\end{itemize}

\begin{tcolorbox}[breakable,colback=teal!4,colframe=teal!35!black,title={Before you move on}]
What does \ja{なんでも} mean? (\textquotedblleft{}Anything; whatever\textquotedblright{}.) Say it once more.
\end{tcolorbox}
